(1) 电磁波在真空中能量传播速度为 \(c\)，动量与能量满足 \(p=E/c\)。若每秒每单位垂直波束的面积传递能量 \(I\)，对应的入射动量通量为 \(I/c\)，其单位 \(\mathrm{W\,m^{-2}}/(\mathrm{m\,s^{-1}})=\mathrm{N\,m^{-2}}\)，即压强。
(2) 完全吸收时，光将每单位面积每秒 \(I/c\) 的法向动量传给表面，因此 \(P_{\rm abs}=I/c\)。同时光场能量完全转化为物体内能，表面吸收热功率密度为 \(I\)。二者分别属于力/面积与能量/(面积·时间)，不应混同。
(3) 无损理想镜使每个光子的法向动量从 \(+E/c\) 变为 \(-E/c\)，改变了 \(-2E/c\)。由动量守恒表面获得相反的 \(+2E/c\)，所以平均辐射压力为 \(P_{\rm mir}=2I/c\)，正好是完全吸收平面的两倍；理想镜不发生净吸收热。
(4) 反射比例 \(\mathcal R\) 带来动量贡献 \(2\mathcal RI/c\)，吸收比例 \(1-\mathcal R\) 贡献 \((1-\mathcal R)I/c\)，合计 \(P=(1+\mathcal R)I/c\)。斜入射角 \(\theta\) 以法线度量时，入射在单位表面积上的能流与法向动量都各含一次 \(\cos\theta\)，理想镜为 \(2I\cos^2\theta/c\)。
